\newpage
\section{Metodologia}
\label{sec:metodologia}

Este capítulo descreve as ferramentas e fontes utilizadas para produzir os
resultados apresentados ao longo do trabalho, organizadas em quatro frentes
metodológicas complementares: revisão bibliográfica estruturada, levantamento do
\textit{market share} das químicas de baterias, construção dos cenários de frota
e Análise de Fluxo de Materiais e Energia (MEFA) de uma fábrica de células LFP.
As três primeiras frentes operam na escala da frota; a quarta, na escala
industrial.

% ===========================================================
\subsection{Revisão Bibliográfica}
\label{sec:met-revisao}
% ===========================================================

A revisão bibliográfica foi conduzida de forma estruturada nas bases de dados
\textbf{Scopus}, \textbf{Web of Science} e \textbf{ScienceDirect}, com cobertura
temporal principal entre 2021 e 2025, de modo a capturar os avanços mais recentes
em um setor de rápida evolução tecnológica. Por não seguir um protocolo
sistemático de seleção (com critérios formais de inclusão e exclusão e
contagem dos documentos triados), a revisão é caracterizada como estruturada.

As buscas foram organizadas em três eixos temáticos:

\begin{enumerate}
    \item \textbf{Química e tecnologia de baterias:} as palavras-chave utilizadas
    incluíram termos como \textit{lithium-ion battery}, \textit{LFP}, \textit{NMC},
    \textit{sodium-ion battery}, \textit{solid-state battery} e \textit{battery
    technology roadmap}. Foram priorizados artigos de revisão com alto número de
    citações, relatórios técnicos de organismos internacionais — como a Agência
    Internacional de Energia (IEA) e o consórcio europeu EUROBAT — e estudos de
    modelagem tecno-econômica que aplicam curvas de aprendizado
    (\textit{Lei de Wright}) para projeção de custos de pacotes de baterias.

    \item \textbf{Processos de reciclagem de baterias de íon-lítio:} as
    palavras-chave empregadas foram \textit{lithium-ion battery recycling},
    \textit{hydrometallurgy spent batteries}, \textit{pyrometallurgy LIB},
    \textit{direct recycling cathode}, \textit{LFP recycling economic},
    \textit{second-life battery} e \textit{battery recycling regulation Brazil}.
    Foram priorizados artigos de revisão abrangentes, estudos de avaliação de ciclo
    de vida e análises tecno-econômicas comparativas, bem como relatórios técnicos
    do Laboratório Nacional de Energia Renovável dos EUA (NREL) e fontes
    regulatórias nacionais.

    \item \textbf{Manufatura de células:} relatórios técnicos e estudos de custo e
    de fluxo de energia da produção de células de íon-lítio
    \cite{knehr2022batpac,degen2022,erakca2021,dai2019}, que fundamentam a MEFA
    descrita na Seção~\ref{sec:met-mefa}.
\end{enumerate}

Para a avaliação da maturidade tecnológica das diferentes químicas e rotas de
reciclagem, adotou-se a escala de \textit{Technology Readiness Level} (TRL),
seguindo a classificação do roteiro EUROBAT para baterias automotivas
(Gerações~3, 4 e~5).

% ===========================================================
\subsection{Levantamento do \textit{Market Share} de Químicas}
\label{sec:met-marketshare}
% ===========================================================

A determinação da participação de mercado das diferentes químicas de baterias na
frota eletrificada brasileira baseou-se no cruzamento de duas fontes de dados
primários:

\begin{itemize}
    \item \textbf{Dados de emplacamentos da ABVE:} registros de vendas de veículos
    eletrificados (BEV, PHEV e HEV) fornecidos pela Associação Brasileira do
    Veículo Elétrico \cite{ABVE}, que constituíram a base quantitativa do estudo.

    \item \textbf{Especificações técnicas do PBEV:} fichas técnicas do Programa
    Brasileiro de Etiquetagem Veicular \cite{PBEV}, utilizadas para identificar
    a química de bateria de cada modelo comercializado no mercado nacional.
\end{itemize}

A metodologia consistiu em (i)~selecionar os modelos de maior impacto em volume de
vendas, segmentados por categoria (BEV, PHEV e HEV); (ii)~atribuir a cada modelo a
respectiva química de bateria com base nas fichas do PBEV e, quando necessário, em
informações técnicas dos fabricantes; e (iii)~consolidar os dados em uma amostra
de \textbf{mais de 509\,mil veículos eletrificados}, a partir da qual foram
calculadas as proporções de cada tecnologia — LFP, NMC e NiMH. O período de
emplacamentos coberto pela base é
\pendente{informar o período da base ABVE (BI Geral) e a data da consulta}.

As participações são calculadas em número de veículos emplacados e, portanto,
não equivalem à participação em massa de baterias ou em capacidade (kWh): o
pacote de um híbrido convencional tem da ordem de 1 a 2~kWh, e o de um BEV,
dezenas de kWh (Seção~\ref{subsec:contexto-brasil}). Essa linha de base
caracteriza a composição química das baterias que circularão e serão
descartadas no país e orienta a escolha da química LFP para o estudo industrial
(Seção~\ref{sec:met-mefa}).

% ===========================================================
\subsection{Construção dos Cenários de Frota}
\label{sec:met-cenarios}
% ===========================================================

Para estabelecer as projeções de eletrificação da frota automotiva brasileira com
horizontes em 2030, 2035 e 2040, foram analisados quatro documentos institucionais
de referência:

\begin{enumerate}
    \item \textbf{ANFAVEA/BCG:} estudo da Associação Nacional dos Fabricantes de
    Veículos Automotores com o \textit{Boston Consulting Group}, divulgado em
    2024, que apresenta dois cenários de eletrificação — transição gradual e
    transição acelerada — para os horizontes de 2030, 2035 e 2040
    \cite{anfaveabcg2024}.

    \item \textbf{PDE 2035 — Plano Decenal de Expansão de Energia}
    \cite{EPE2025}: projeções governamentais de vendas e frota circulante de
    veículos leves e pesados, com detalhamento do eixo de eletromobilidade.

    \item \textbf{Bright Consulting:} previsões de mercado e frota com horizonte
    principal em 2030, com foco na trajetória de hibridização da frota brasileira
    \cite{BrightConsulting}.

    \item \textbf{MBC Brasil — Movimento Brasil Competitivo:} estudo sobre
    transição energética da mobilidade de baixo carbono no Brasil com horizonte em
    2040 \cite{MBCBrasil}.
\end{enumerate}

Diante da disparidade numérica entre as projeções, adotou-se como espinha
dorsal o estudo ANFAVEA/BCG — única fonte consultada que apresenta dois cenários
e os três horizontes de interesse — e as demais fontes foram usadas como
verificação cruzada de ordem de grandeza:

\begin{itemize}
    \item \textbf{Cenário Conservador:} corresponde ao cenário de transição gradual
    do estudo ANFAVEA/BCG, com predominância do Híbrido-Flex e gargalos de
    infraestrutura de recarga.

    \item \textbf{Cenário Acelerado:} corresponde ao cenário de transição
    acelerada, com maior penetração de BEV e PHEV, associada à efetividade de
    programas como o \textit{Mover} e o \textit{Novo PAC}.
\end{itemize}

Os valores por categoria e horizonte são reproduzidos da fonte, arredondados ao
milhão (Seção~\ref{subsec:cenarios-sensibilidade}). As projeções do PDE 2035, da
Bright Consulting e do MBC Brasil são comparadas a eles no mesmo horizonte,
distinguindo-se, quando possível, a inclusão ou não dos híbridos leves (MHEV).

Os veículos pesados não integram a quantificação dos cenários. Para os
caminhões, as projeções consultadas indicam predominância do diesel ao longo do
horizonte (Seção~\ref{subsec:pesados}); para os ônibus, embora a eletrificação
urbana seja relevante, as fontes consultadas não oferecem séries comparáveis às
dos veículos leves para os três horizontes. Os cenários quantitativos
concentram-se, portanto, nos \textbf{veículos leves}.

% ===========================================================
\subsection{Análise de Fluxo de Materiais e Energia da Fábrica}
\label{sec:met-mefa}
% ===========================================================

A Análise de Fluxo de Materiais (MFA) é um sistema de equações de balanço de
massa aplicado a processos e fluxos delimitados no espaço e no tempo
\cite{BrunnerRechberger}. Neste trabalho, o método é aplicado a uma fábrica de
células LFP e estendido aos fluxos de energia e de utilidades — eletricidade,
água e gases técnicos —, à semelhança da análise de fluxos de energia aplicada
por \citeonline{erakca2021} à produção de células; por isso, adota-se a
denominação Análise de Fluxo de Materiais e Energia (MEFA). Os elementos do
método são:

\begin{itemize}
    \item \textbf{Unidade funcional:} 1~kWh de célula LFP acabada e testada, em
    formato prismático, com densidade energética de referência de 160~Wh/kg
    (6,25~kg de célula por kWh).

    \item \textbf{Base de normalização:} os resultados por kWh são multiplicados
    por $10^{6}$ e expressos por GWh de células produzidas por ano. A
    normalização não representa a capacidade física de uma planta específica: os
    parâmetros correspondem a uma planta madura em operação consolidada, e plantas
    pequenas ou em \textit{ramp-up} apresentam intensidades energéticas superiores
    (Seção~\ref{subsec:giga-indicadores}).

    \item \textbf{Fronteira:} do tipo \textit{gate-to-gate}, do recebimento das
    matérias-primas no portão da fábrica até a expedição das células acabadas; as
    inclusões e exclusões são detalhadas na Seção~\ref{subsec:giga-bases}.

    \item \textbf{Fontes:} faixas de referência da literatura técnica de
    manufatura de células \cite{knehr2022batpac,degen2022,erakca2021,dai2019};
    não há medições de uma planta específica.

    \item \textbf{Tratamento das faixas:} cada parâmetro é registrado com a faixa
    da literatura e um valor central adotado para o balanço. Os valores são
    classificados como \textbf{L} (dado de literatura), \textbf{H} (hipótese
    adotada neste trabalho) ou \textbf{C} (calculado a partir de L e H),
    classificação indicada nas tabelas do Capítulo~\ref{sec:gigafabrica}.

    \item \textbf{Balanços:} em regime permanente, para cada processo $i$,
    \begin{equation}
        \sum_{j} \dot{m}^{\,\text{entrada}}_{i,j} = \sum_{k} \dot{m}^{\,\text{saída}}_{i,k},
        \label{eq:balanco}
    \end{equation}
    e o balanço global resulta da composição dos balanços por etapa, aplicados aos
    materiais que efetivamente atravessam cada etapa.

    \item \textbf{Classificação dos fluxos:} (a) entradas externas, que cruzam a
    fronteira vindas de fornecedores; (b) saídas externas — produto, subprodutos,
    resíduos, emissões e efluentes; e (c) circuitos internos, que deixam uma
    etapa e retornam ao processo sem cruzar a fronteira (recuperação de NMP,
    retorno de pasta, recirculação de água e de energia).
\end{itemize}

% ===========================================================
\subsection{Síntese da Abordagem Metodológica}
\label{sec:met-sintese}
% ===========================================================

A Figura~\ref{fig:fluxo-metodologico} ilustra a integração entre as quatro
frentes. Na escala da frota, os dados de emplacamento (ABVE/PBEV) alimentam o
\textit{market share}, e os documentos institucionais alimentam os cenários. Na
escala industrial, a literatura de manufatura alimenta a MEFA da fábrica, cuja
química decorre do \textit{market share}. As duas escalas convergem, com o
contexto tecnológico e regulatório da revisão bibliográfica, para a discussão das
implicações para a cadeia de reciclagem. A modelagem dinâmica dos fluxos da frota (MFA em uso, fim de vida
e segunda vida) é indicada como etapa futura.

\begin{figure}[H]
\centering
\resizebox{\textwidth}{!}{%
\begin{tikzpicture}[
    box/.style   = {rectangle, rounded corners, draw=black!60, thick,
                    fill=gray!8, align=flush center, minimum height=1.1cm,
                    text width=3.3cm, font=\footnotesize},
    fonte/.style = {box, fill=blue!6},
    res/.style   = {box, fill=orange!10},
    ind/.style   = {box, fill=green!10},
    sint/.style  = {box, fill=purple!8, text width=3.0cm, minimum height=2.6cm},
    fut/.style   = {box, dashed, fill=white, draw=black!50, text=black!70,
                    text width=3.0cm},
    arrow/.style = {-{Latex[length=2.6mm]}, thick},
    darrow/.style= {-{Latex[length=2.6mm]}, thick, dashed, black!55},
]
% --- escala da frota
\node[fonte] (inst) at (0, 0)    {Documentos institucionais\\(ANFAVEA/BCG, PDE, Bright, MBC)};
\node[fonte] (abve) at (0,-1.9)  {ABVE / PBEV\\(emplacamentos\\e fichas)};
\node[res]   (cen)  at (5, 0)    {Cenários de frota\\(2030, 2035, 2040)};
\node[res]   (ms)   at (5,-1.9)  {\textit{Market share}\\das químicas};
% --- escala industrial
\node[fonte] (man)  at (0,-4.2)  {Literatura de manufatura\\de células};
\node[ind]   (mefa) at (5,-4.2)  {MEFA da fábrica\\de células LFP};
% --- revisão
\node[fonte] (lit)  at (0,-6.1)  {Literatura científica\\e técnica};
\node[box]   (rev)  at (5,-6.1)  {Revisão\\bibliográfica};
% --- molduras de escala
\begin{scope}[on background layer]
\node[draw=black!30, dashed, rounded corners, inner sep=7pt, fill=gray!3,
      fit=(inst)(abve)(cen)(ms),
      label={[font=\scriptsize\bfseries, anchor=south west]north west:Escala da frota}]
      (frota) {};
\node[draw=black!30, dashed, rounded corners, inner sep=7pt, fill=gray!3,
      fit=(man)(mefa),
      label={[font=\scriptsize\bfseries, anchor=south west]north west:Escala industrial}]
      (industria) {};
\end{scope}
% --- síntese e etapa futura
\node[sint] (impl) at (10.3,-3.05) {Implicações\\para a cadeia\\de reciclagem};
\node[fut]  (mfa)  at (10.3, 0.35) {MFA dinâmico\\da frota\\(etapa futura)};
% --- setas
\draw[arrow] (inst) -- (cen);
\draw[arrow] (abve) -- (ms);
\draw[arrow] (man)  -- (mefa);
\draw[arrow] (lit)  -- (rev);
\draw[arrow] (ms) -- node[right=2pt, font=\scriptsize, align=left] {química LFP} (mefa);
\draw[arrow] (frota.east |- 0,-2.0) -- (impl.west |- 0,-2.0);
\draw[arrow] (industria.east |- mefa.east) -- (impl.west |- mefa.east);
\draw[arrow] (rev.east) -| (impl.south);
\draw[darrow] (frota.east |- mfa.west) -- (mfa.west);
\end{tikzpicture}%
}
\caption{Fluxo metodológico integrado da pesquisa. Setas tracejadas indicam a
         etapa futura (MFA dinâmico da frota), fora do escopo deste relatório.}
\label{fig:fluxo-metodologico}
\end{figure}
